\section{Harness and Method}
\label{sec:method}

\subsection{One interface, interchangeable deciders}

A decider receives a rendered state, an instruction, and a set of options, and returns one option with a confidence.
Every decider at a decision point sees the identical rendered state and the identical option descriptions.
We evaluate the following deciders.

\textbf{Decision models.} \jev\ 1.13 through its hosted API; \kev\ (4B by default, also 0.8B and 9B for agent control) served locally.

\textbf{LLMs.} \gbig\ as the frontier ceiling; \goss\ as a small hosted LLM; \qwen~\cite{yang2025qwen3} served locally with vLLM~\cite{kwon2023vllm}, with thinking turned off.
An LLM decides through one constrained call: the options are the only permitted outputs (a forced tool choice), and its confidence is the one it states.

\textbf{Baselines specific to a decision point.}
For the cache guard: a cut-off on the embedding similarity of the two queries; an off-the-shelf duplicate-question cross-encoder~\cite{reimers2019sbert}; vCache~\cite{schroeder2025vcache}, which learns a similarity cut-off per cached prompt under a target error rate $\delta$; and a classifier fine-tuned on each dataset's dev split.
For the router: a classifier fine-tuned on the cheap model's record of right and wrong answers, plus three fixed policies (always small, always large, and a random mix of the two).

\textbf{Cascades.} A cheap first stage answers when it is confident and escalates otherwise: \jev\ to \gbig\ (agent control), and similarity threshold to \jev\ (cache guard).
Every cascade threshold or band is chosen on data the reported number does not use.

\subsection{Measurement}

\textbf{Repeats and intervals.} Stochastic deciders are called five times per item (\gbig\ three times at the cache guard and router, to bound cost).
Repeats on one item are correlated, so all intervals are 95\% cluster bootstraps: over states or tasks for agent control, over classes for the cache guard, over prompts for the router.
Comparisons between two deciders are paired on the same items.

\textbf{Cost.} For hosted models we record the cost the API gateway reports for each call, not a token count multiplied by a list price; the two differed by about $3\times$ for one model because of upstream provider selection.
All hosted calls go through one gateway and one key.
Calls from different deciders are interleaved so that network conditions affect them equally.

\textbf{Latency.} Wall-clock time per call at the client.
For local deciders we also record server-side model time, which isolates HTTP overhead (7--8\,ms).

\textbf{Energy.} For local deciders we read the GPU's cumulative energy counter (NVML) around a contiguous block of calls on a dedicated A100-80GB, with one request in flight, prefix caching off, and warm-up calls discarded.
We report gross joules per decision and the value net of the idle baseline.
The counter is validated against the instantaneous power reading before each run; on one consumer GPU it failed that check, and we report no energy for that machine.
Hosted deciders have no energy number.

\textbf{Failures.} A call that fails for an infrastructure reason (billing, rate limit, provider error, network) is retried and never counted against the decider.
A call in which the decider produces no valid choice is a decider failure, counted as wrong and never retried, so a stochastic decider cannot re-roll its own failures.

\textbf{Calibration.} Expected calibration error (ECE)~\cite{guo2017calibration} over the probability of the chosen option.
